(1) 球形晶核面积为 \(4\pi r^2\)，形成界面需要自由能 \(4\pi r^2\gamma\)；转化为新相的体积为 \(4\pi r^3/3\)，有利于相变的体积自由能降低 \(4\pi r^3\Delta g_v/3\)。两者相减得 \(\Delta G=4\pi r^2\gamma-4\pi r^3\Delta g_v/3\)。
(2) 求导得 \(\Delta G'=4\pi r(2\gamma-\Delta g_vr)\)。非零驻点是 \(r_*=2\gamma/\Delta g_v\)，代入原自由能函数得 \(\Delta G_*=(16\pi/3)\gamma^3/\Delta g_v^2\)，这就是经典均匀成核势垒。
(3) 二阶导数 \(\Delta G''=8\pi\gamma-8\pi\Delta g_vr\)，在临界点 \(\Delta G''(r_*)=-8\pi\gamma<0\)，因此临界晶核是势垒顶端而非稳定尺寸。比 \(r_*\) 小的晶核倾向收缩，比其大的核可沿增大半径方向降低自由能。
(4) 临界半径按 \(r_*\propto\gamma/\Delta g_v\) 缩放，势垒按 \(\Delta G_*\propto\gamma^3/\Delta g_v^2\) 缩放。因此界面张力提高会抑制成核，而增大过冷导致的体积驱动力会促进成核；经典近似要求足够大的宏观核，且不考虑界面各向异性。
